\documentclass[11pt,a4paper]{article}
\usepackage[utf8]{inputenc}
\usepackage[T1]{fontenc}
\usepackage{amsmath,amssymb,amsfonts}
\usepackage{booktabs}
\usepackage{graphicx}
\usepackage{multirow}
\usepackage{hyperref}
\usepackage[margin=2.5cm]{geometry}
\usepackage{xcolor}
\usepackage{enumitem}

\title{Compounding Failures in Per-Head KV Cache Quantization:\\An Empirical Analysis of When Non-Uniform Allocation Helps}

\author{Anonymous}
\date{}

\begin{document}
\maketitle

%% ============================================================
%% ABSTRACT
%% ============================================================
\begin{abstract}
Per-head mixed-precision KV cache quantization promises better quality--compression tradeoffs than uniform bit-width assignment.
We show that on hardware-aligned non-consecutive bit-widths $\{2,4,8,16\}$, this promise breaks:
2-bit quantization distortions compound non-additively through the residual stream, causing \emph{all} tested unconstrained solvers (greedy, MCKP) to over-allocate 2-bit heads and underperform uniform 4-bit across 8 diverse architectures (degradation $1.07\times$--$3.31\times$ baseline).
A simple cardinality constraint ($\le k$ heads at 2-bit) prevents compounding across all tested architectures, and reduces the feasible set to where solver choice is irrelevant: constrained greedy and constrained MCKP produce identical bit allocations ($\Delta=0.0000$ PPL).
Surprisingly, under moderate (4$\times$) compression, only 1 of 8 tested architectures benefits from non-uniform allocation; the remaining 7 show no benefit or minor degradation.
We propose a 2-metric empirical diagnostic (top-5\% concentration $\gtrsim 0.50$ AND $t_1/t_5 \lesssim 0.70$) that correctly classifies all 9 tested models, giving practitioners a simple check: if the diagnostic says ``no,'' skip profiling and use uniform 4-bit with zero quality loss.
At aggressive compression (8$\times$), the constraint reduces degradation from $45.9\times$ to $5.9\times$ baseline on Qwen2.5-7B, demonstrating critical deployment value.
\end{abstract}

%% ============================================================
%% 1  INTRODUCTION
%% ============================================================
\section{Introduction}
\label{sec:intro}

KV cache memory is a primary bottleneck for long-context LLM serving.
Per-head mixed-precision quantization---assigning different bit-widths to different attention heads based on their sensitivity---has been proposed as a principled approach to improve the quality--compression tradeoff over uniform quantization~\cite{hooper2024kvquant,liu2024kivi,dong2024qaq}.
The underlying assumption is straightforward: allocate fewer bits to insensitive heads and more bits to sensitive ones.

However, modern hardware efficiently supports only non-consecutive bit-widths $\{2,4,8,16\}$, creating a large gap between the lowest two options.
We discover that this gap introduces a previously unreported failure mode: \textbf{2-bit quantization distortions compound non-additively} when multiple heads are simultaneously assigned 2-bit precision.
Per-head profiling predicts a joint PPL of ${\sim}11$ when assigning many heads to 2-bit, but the actual PPL is 32--37---a $3\times$ prediction error.
This compounding causes all tested unconstrained allocation algorithms to systematically over-allocate 2-bit heads, producing worse results than naive uniform 4-bit across \emph{all 8 tested architectures} (degradation ranging from $1.07\times$ to $3.31\times$ baseline).

This paper makes three contributions, framed as an empirical analysis:

\begin{enumerate}[leftmargin=*,itemsep=2pt]
\item \textbf{Compounding effect characterization.} We provide the first quantitative mapping of 2-bit compounding in per-head KV cache quantization, showing superlinear PPL growth after $k \approx 16$ simultaneous 2-bit heads and exponential growth after $k \approx 48$ on Qwen2.5-7B (Table~\ref{tab:compounding_sweep}).
The pattern is consistent across 8 architectures spanning GQA ratios 2:1 to 8:1.

\item \textbf{Cardinality constraint as universal fix.} A cardinality constraint ($\le k$ heads at 2-bit) prevents compounding across all tested architectures.
Under this constraint, solver choice becomes irrelevant: constrained greedy and constrained MCKP produce \emph{identical} allocations ($\Delta = 0.0000$; Table~\ref{tab:solver_equivalence}).
The safe range $k \le 8$ yields $<2.5\%$ degradation across all tested architectures.

\item \textbf{When does non-uniform allocation help?} Under moderate ($4\times$) compression, only 1/8 architectures (Qwen2.5-7B, ``spread'' sensitivity pattern) benefits from non-uniform allocation (+27.3\% PPL improvement, +2.6pp average downstream gain).
The remaining 7/8 show no benefit or minor degradation.
A 2-metric diagnostic (9/9 correct) tells practitioners whether profiling is worthwhile.
\end{enumerate}

Our work differs from RateQuant~\cite{ratequant2024} in three key aspects: (a)~RateQuant operates on consecutive bit-widths $\{2,\ldots,6\}$ where polymatroid structure guarantees greedy optimality; on hardware-aligned $\{2,4,8,16\}$ this structure breaks, (b)~RateQuant acknowledges compounding in Appendix~M but does not characterize it or propose a fix---we provide the first quantitative sweep and a universal constraint, (c)~we provide the first systematic study of \emph{when} non-uniform allocation is actually beneficial versus when uniform suffices.


%% ============================================================
%% 2  RELATED WORK
%% ============================================================
\section{Related Work}
\label{sec:related}

\paragraph{KV Cache Quantization.}
KIVI~\cite{liu2024kivi} applies uniform per-channel quantization to KV caches with asymmetric K/V treatment.
KVQuant~\cite{hooper2024kvquant} introduces per-channel sensitivity-aware quantization with non-uniform codebooks.
QAQ~\cite{dong2024qaq} proposes attention-aware quantization with quality-adaptive bit-width selection.

\paragraph{Mixed-Precision Allocation.}
RateQuant~\cite{ratequant2024} formulates per-head bit allocation as a rate--distortion problem, deriving a continuous closed-form solution (Thm.~2) and a discrete greedy algorithm (Prop.~6) that is provably optimal on consecutive integer bit-widths.
On non-consecutive hardware-aligned widths $\{2,4,8,16\}$, the polymatroid structure underlying this optimality guarantee breaks.
RateQuant acknowledges compounding effects in Appendix~M but provides no quantitative characterization or mitigation.
Our work provides the first such characterization and shows that a cardinality constraint both fixes the compounding failure and makes the solver choice irrelevant---constrained greedy matches constrained MCKP exactly.

\paragraph{Compounding Effects in Quantization.}
Non-additive error accumulation is well-known in weight quantization~\cite{nagel2020up}, but its manifestation in per-head KV cache quantization---where the residual stream propagates head-level errors across layers---has not been systematically studied.
We provide the first quantitative evidence that per-head independent profiling underestimates joint 2-bit degradation by ${\sim}3\times$.


%% ============================================================
%% 3  METHOD
%% ============================================================
\section{Method}
\label{sec:method}

\subsection{Problem Formulation}

Given an LLM with $H$ KV heads and a set of hardware-aligned bit-widths $\mathcal{B} = \{2, 4, 8, 16\}$, per-head mixed-precision allocation assigns a bit-width $b_h \in \mathcal{B}$ to each head $h$ subject to a memory budget constraint:
\begin{equation}
\min_{\{b_h\}} \; \text{PPL}(\{b_h\}) \quad \text{s.t.} \quad \sum_{h=1}^{H} b_h \le B
\label{eq:allocation}
\end{equation}
where $B$ is the total bit budget.

\subsection{Per-Head Sensitivity Profiling}

For each head $h$ and candidate bit-width $b \in \mathcal{B}$, we measure the \emph{distortion delta}: the PPL increase when head $h$ alone is quantized to $b$-bit while all other heads remain at 16-bit.
This produces a distortion matrix $\Delta(h, b)$ of size $H \times |\mathcal{B}|$.

Profiling uses calibration data from WikiText-2 (train split).
We evaluate two resolutions: low-resolution (16 sequences $\times$ 512 tokens) and high-resolution (64 sequences $\times$ 2048 tokens).
High-resolution profiling detects sensitivity in 100/112 heads versus 37/112 for low-resolution on Qwen2.5-7B, improving constrained PPL from 11.23 to 10.91 ($-2.8\%$; Table~\ref{tab:profiling_resolution}).

\subsection{MCKP Formulation with Cardinality Constraint}

We formulate bit-width allocation as a Multiple-Choice Knapsack Problem (MCKP): each head selects exactly one bit-width from $\mathcal{B}$, with the objective of minimizing total distortion subject to the memory budget.
Critically, we add a \textbf{cardinality constraint}:
\begin{equation}
\sum_{h=1}^{H} \mathbf{1}[b_h = 2] \le k
\label{eq:constraint}
\end{equation}
limiting the number of heads assigned to 2-bit quantization.
This constraint is motivated by our empirical observation that 2-bit distortions compound non-additively (Section~\ref{sec:compounding}).

The constraint parameter $k$ is selected via a small-scale compounding test: sweep $k$ on a validation set and identify the ``safe'' range where PPL growth remains approximately linear.
Across all tested architectures, $k \le 8$ consistently yields $<2.5\%$ degradation relative to the unconstrained optimal (Table~\ref{tab:compounding_sweep}).

\subsection{Applicability Diagnostic}

Not all architectures benefit from non-uniform allocation.
We propose a 2-metric diagnostic based on the per-head sensitivity profile:

\begin{itemize}[leftmargin=*,itemsep=2pt]
\item \textbf{Top-5\% concentration} ($c$): the fraction of total distortion delta concentrated in the top 5\% of heads.
\item \textbf{$t_1/t_5$ ratio}: the ratio of the largest single-head distortion to the sum of the top-5 distortions.
\end{itemize}

The decision rule: if $c \gtrsim 0.50$ AND $t_1/t_5 \lesssim 0.70$, the model has a ``spread'' sensitivity pattern and non-uniform allocation is recommended.
Otherwise, uniform 4-bit suffices at moderate compression.
This diagnostic correctly classifies all 9 tested models (Section~\ref{sec:when_helps}).


%% ============================================================
%% 4  COMPOUNDING EFFECT CHARACTERIZATION
%% ============================================================
\section{Compounding Effect Characterization}
\label{sec:compounding}

\subsection{Compounding Sweep}

Table~\ref{tab:compounding_sweep} shows the PPL of Qwen2.5-7B as a function of $k$ (maximum 2-bit heads) under $4\times$ compression (budget$=0.25$).
PPL growth is approximately linear for $k \le 8$ ($+2.4\%$), transitions to superlinear at $k \approx 16$ ($+6.2\%$), and becomes exponential after $k \approx 48$ ($+136\%$), reaching $10.07\times$ baseline at $k = 112$ (fully unconstrained).

\begin{table}[h]
\centering
\caption{Compounding Sweep on Qwen2.5-7B (budget$=0.25$). PPL grows superlinearly with the number of simultaneous 2-bit heads.}
\label{tab:compounding_sweep}
\begin{tabular}{@{}rrrr@{}}
\toprule
$k$ (max 2-bit heads) & PPL & Ratio to $k{=}0$ \\
\midrule
0  & 9.81  & 1.00$\times$ \\
8  & 10.05 & 1.02$\times$ \\
16 & 10.42 & 1.06$\times$ \\
32 & 11.60 & 1.18$\times$ \\
64 & 23.22 & 2.37$\times$ \\
112 & 98.75 & 10.07$\times$ \\
\bottomrule
\end{tabular}
\end{table}

\subsection{Cross-Architecture Unconstrained Degradation}

The compounding failure mode is consistent across all 8 tested architectures.
Without a cardinality constraint, solvers over-allocate 2-bit heads and \emph{always} underperform uniform 4-bit.
The degradation follows a concentration-dependent gradient:
\begin{itemize}[leftmargin=*,itemsep=1pt]
\item \textbf{Catastrophic} ($3.31\times$): Qwen2.5-7B (spread pattern, L0 cluster compounding)
\item \textbf{Moderate} ($1.47\times$): Gemma-2-9B (flat pattern, solver over-allocates 106/336 heads to 2-bit due to indistinguishable distortions)
\item \textbf{Mild} ($1.07$--$1.08\times$): Yi-1.5-9B, Llama-3.1-8B, Mistral-7B-v0.3
\end{itemize}

Gemma-2-9B is instructive: despite its flat sensitivity distribution (concentration$=0.23$, CV$=1.28$), the solver assigns 31.5\% of heads to 2-bit because distortions are indistinguishable.
The cumulative effect of this over-allocation produces a moderate $1.47\times$ degradation.
The cardinality constraint reduces 2-bit heads from 106 to 8, yielding PPL within 0.3\% of uniform---confirming that the constraint universally prevents solver over-allocation regardless of distribution shape.


%% ============================================================
%% 5  EXPERIMENTS
%% ============================================================
\section{Experiments}
\label{sec:experiments}

All experiments use WikiText-2 for calibration and evaluation unless otherwise noted.
We report perplexity (PPL) on the test set with a context length of 2048 tokens and $4\times$ compression (budget$=0.25$) as the default setting.

\subsection{Cross-Architecture Pipeline Results}

Table~\ref{tab:cross_arch} summarizes the full pipeline results across three primary architectures.
The constrained MCKP benefit varies dramatically: Qwen2.5-7B shows a large improvement (+27.3\% vs uniform), Llama-3.1-8B shows marginal improvement (+0.2\%), and Mistral-7B-v0.3 shows degradation ($-0.5\%$).

\begin{table}[h]
\centering
\caption{Cross-architecture pipeline results (budget$=0.25$). Constrained MCKP benefits are architecture-dependent.}
\label{tab:cross_arch}
\begin{tabular}{@{}lrrrrr@{}}
\toprule
Model & BF16 & Uniform 4-bit & Constrained MCKP & $\Delta$ vs Uniform \\
\midrule
Qwen2.5-7B & 9.81 & 15.45 & 10.91 & +27.3\% \\
Llama-3.1-8B & 8.09 & 8.47 & 8.45 & +0.2\% \\
Mistral-7B-v0.3 & 7.04 & 7.13 & 7.17 & $-$0.5\% \\
\bottomrule
\end{tabular}
\end{table}

\paragraph{Llama-3.1-8B.}
The solver naturally selects only 2 heads for 2-bit (distribution $\{2{:}2, 4{:}253, 8{:}1\}$), so the cardinality constraint ($k{=}8$) never binds.
The +0.2\% improvement over uniform is within noise.
This is consistent with Llama's spike sensitivity distribution ($t_1/t_5 \approx 0.85$): a single dominant head (L2H5) provides no optimization opportunity for the solver.

\paragraph{Mistral-7B-v0.3.}
PPL monotonically increases with $k$: every additional 2-bit head hurts.
Even the best constrained setting ($k{=}4$, PPL$=7.165$) is worse than uniform 4-bit ($7.131$).
This reflects Mistral's flat sensitivity distribution (concentration$=0.22$): no head benefits from 2-bit downgrade.

\begin{table}[h]
\centering
\caption{Mistral-7B $k$-ablation (budget$=0.25$). PPL increases monotonically with $k$; uniform 4-bit is optimal.}
\label{tab:mistral_k_ablation}
\begin{tabular}{@{}rrrl@{}}
\toprule
$k$ & PPL & \#2-bit heads & Note \\
\midrule
Uniform & 7.131 & 0 & Optimal for Mistral \\
4  & 7.165 & 4  & Best constrained \\
8  & 7.200 & 8  & \\
16 & 7.311 & 16 & \\
\bottomrule
\end{tabular}
\end{table}

\subsection{Solver Equivalence Under Constraint}

Table~\ref{tab:solver_equivalence} presents the strongest evidence that the cardinality constraint is the primary contribution.
Without the constraint, both greedy and MCKP fail catastrophically (PPL$>$30).
With the constraint ($k{=}8$), both produce \emph{identical} bit allocations ($\{2{:}8, 4{:}100, 8{:}4\}$) and PPL ($\Delta = 0.0012$).

\begin{table}[h]
\centering
\caption{Solver equivalence under cardinality constraint (Qwen2.5-7B, $k{=}8$). Both solvers fail unconstrained; with the constraint, they produce identical allocations.}
\label{tab:solver_equivalence}
\begin{tabular}{@{}lrll@{}}
\toprule
Solver & PPL & Bit Distribution & Status \\
\midrule
Greedy unconstrained & 32.45 & $\{2{:}68, 4{:}30, 8{:}4, 16{:}10\}$ & FAIL \\
MCKP unconstrained & 36.92 & --- & FAIL \\
Greedy + max\_2bit$=$8 & 11.23 & $\{2{:}8, 4{:}100, 8{:}4\}$ & PASS \\
MCKP + max\_2bit$=$8 & 11.23 & $\{2{:}8, 4{:}100, 8{:}4\}$ & PASS \\
Uniform 4-bit & 15.45 & all 4-bit & baseline \\
\bottomrule
\end{tabular}
\end{table}

This equivalence demonstrates that the constraint reduces the feasible region to where solver sophistication is irrelevant.
The practical implication is that practitioners can add a single integer parameter (\texttt{max\_2bit}) to \emph{any} existing allocation pipeline---including RateQuant's greedy---and achieve near-optimal results.

\subsection{MCKP Near-Optimality (Oracle Comparison)}

To quantify how close MCKP is to the true optimum within the constrained feasible set, we perform brute-force enumeration over 4-layer groups on Qwen2.5-7B (Table~\ref{tab:oracle}).
MCKP is exactly optimal in 5/8 layer groups; the remaining 3 show gaps of 0.0045--0.0136 PPL, arising from 2-bit head position swaps within the same bit budget.
The maximum oracle gap is 0.018 PPL---negligible compared to the 4.54 PPL gap between constrained MCKP (10.91) and uniform 4-bit (15.45).

\begin{table}[h]
\centering
\caption{MCKP vs brute-force oracle per layer group (Qwen2.5-7B). MCKP is near-optimal within the constrained feasible set.}
\label{tab:oracle}
\begin{tabular}{@{}lcr@{}}
\toprule
Layer Group & MCKP Optimal? & Gap (PPL) \\
\midrule
L0--L3  & Yes & 0.000 \\
L4--L7  & No  & 0.005 \\
L8--L11 & No  & 0.014 \\
L12--L15 & Yes & 0.000 \\
\bottomrule
\end{tabular}
\end{table}

\subsection{Budget Sweep (Compression Ratio)}

Table~\ref{tab:budget_sweep} shows how the constrained MCKP benefit scales with compression ratio on Qwen2.5-7B.
At $2\times$ compression, uniform and MCKP are equivalent ($\sim 0\%$ difference).
At $4\times$, MCKP provides 29.4\% improvement.
At $8\times$, the gap becomes dramatic: MCKP achieves PPL$=57.6$ versus uniform 2-bit PPL$=2643.77$ ($45.9\times$)---the constraint is the difference between a usable and unusable model at aggressive compression.

\begin{table}[h]
\centering
\caption{Budget sweep on Qwen2.5-7B. Compression amplifies the benefit of constrained allocation.}
\label{tab:budget_sweep}
\begin{tabular}{@{}lrrl@{}}
\toprule
Compression & MCKP PPL & Uniform PPL & $\Delta$ \\
\midrule
$2\times$ (budget$=0.5$) & 9.85 & 9.86 & $\sim 0\%$ \\
$4\times$ (budget$=0.25$) & 10.91 & 15.45 & 29.4\% \\
$8\times$ (budget$=0.125$) & 57.60 & 2643.77 & 45.9$\times$ \\
\bottomrule
\end{tabular}
\end{table}

\subsection{Profiling Resolution Impact}

\begin{table}[h]
\centering
\caption{Profiling resolution impact (Qwen2.5-7B, $k{=}8$). Higher resolution improves allocation quality at increased profiling cost.}
\label{tab:profiling_resolution}
\begin{tabular}{@{}lrrl@{}}
\toprule
Resolution & PPL & Detected Heads & Time (min) \\
\midrule
Low-res ($16 \times 512$)  & 11.23 & 37/112  & $\sim$10 \\
High-res ($64 \times 2048$) & 10.91 & 100/112 & 65.5 \\
Uniform 4-bit & 15.45 & N/A & 0 \\
\bottomrule
\end{tabular}
\end{table}

High-resolution profiling reduces PPL from 11.23 to 10.91 ($-2.8\%$), detecting sensitivity in 100/112 heads versus 37/112 for low-resolution.
However, even low-resolution profiling yields a 27.3\% improvement over uniform 4-bit (11.23 vs 15.45), suggesting that approximate profiling suffices for the primary benefit.
The 65.5-minute profiling cost at high resolution is a practical consideration; we note this as an engineering tradeoff rather than a fundamental limitation.

\subsection{L0 Exclusion Ablation}

On Qwen2.5-7B, Layer 0 heads exhibit a $26\times$ dominance ratio in sensitivity delta.
We test whether explicit L0 exclusion (hard-coding L0 heads to 8-bit) is necessary, or whether the solver discovers L0 sensitivity automatically from profiling data.

\begin{table}[h]
\centering
\caption{L0 exclusion ablation (Qwen2.5-7B, budget$=0.25$). The solver automatically assigns 8-bit to all 4 L0 heads.}
\label{tab:l0_ablation}
\begin{tabular}{@{}lrll@{}}
\toprule
Condition & PPL & L0 Bit-width & Note \\
\midrule
MCKP + L0 exclusion + $k{=}8$ & 10.91 & 8-bit (hard-coded) & Original \\
MCKP + $k{=}8$ only & 10.91 & 8-bit (solver-discovered) & Automatic \\
\bottomrule
\end{tabular}
\end{table}

The solver automatically assigns 8-bit to all 4 L0 heads, producing identical PPL.
This confirms that the profiling-based approach correctly identifies L0 sensitivity without manual intervention, though we note this has only been verified on Qwen2.5-7B where L0 dominance is extreme ($26\times$).

\subsection{K/V Separation}

We test whether treating Key and Value heads as independent MCKP items with separate budgets improves over joint K$=$V treatment (Table~\ref{tab:kv_sep}).

\begin{table}[h]
\centering
\caption{K/V separation vs joint treatment (Qwen2.5-7B, $k{=}8$). Separation is worse due to K-cache compounding.}
\label{tab:kv_sep}
\begin{tabular}{@{}lrrr@{}}
\toprule
Method & PPL & K 2-bit & V 2-bit \\
\midrule
Joint K$=$V ($k{=}8$) & 10.91 & 8 & 0 \\
K/V Separation & 11.22 & 16 & 0 \\
Uniform 4-bit & 15.45 & 0 & 0 \\
\bottomrule
\end{tabular}
\end{table}

K/V separation is worse ($+0.31$ PPL) because it allows 16 K-heads at 2-bit, exceeding the safe range of $k \le 8$ and triggering K-cache compounding.
V-heads are assigned uniform 4-bit under both settings.
This suggests that K-cache is more sensitive to 2-bit compounding than V-cache, and that a single joint cardinality constraint is preferable.

\subsection{Heuristic Baselines}

We compare constrained MCKP against simpler allocation strategies: L0-random (protect L0, randomly assign remaining 2-bit heads) and naive positional (assign 2-bit by layer position).

\begin{table}[h]
\centering
\caption{MCKP vs heuristic baselines (budget$=0.25$). MCKP consistently outperforms random constraint-respecting allocation.}
\label{tab:heuristic_baselines}
\begin{tabular}{@{}lrrrr@{}}
\toprule
Model & MCKP & L0-random (median) & Naive Positional & Uniform \\
\midrule
Qwen2.5-7B  & 10.91 & 11.52 & 12.23 & 15.45 \\
Llama-3.1-8B & 8.45 & 10.04 & 9.94 & 8.47 \\
Mistral-7B-v0.3 & 7.16 & 7.59 & 7.48 & 7.13 \\
\bottomrule
\end{tabular}
\end{table}

MCKP outperforms L0-random by 5.2\% (Qwen), 15.8\% (Llama), and 5.7\% (Mistral) in PPL, demonstrating that profiling-based allocation provides optimization value within the constrained feasible set.
The large seed variance in L0-random (Qwen std$=0.47$, Llama std$=1.03$) indicates that the specific choice of 2-bit heads matters significantly.

\subsection{Cross-Dataset Calibration Robustness}

\begin{table}[h]
\centering
\caption{Cross-dataset calibration agreement ($>$96\% across 3 architectures). Profiling-based sensitivity ranking is robust to calibration source.}
\label{tab:cross_dataset}
\begin{tabular}{@{}lrrl@{}}
\toprule
Model & Agreement & Same-domain PPL & Cross-domain PPL \\
\midrule
Qwen2.5-7B  & 108/112 (96.4\%) & 11.29 & 20.56 \\
Llama-3.1-8B & 250/256 (97.7\%) & 8.45 & 8.65 \\
Mistral-7B-v0.3 & 247/256 (96.5\%) & 7.16 & 7.19 \\
\bottomrule
\end{tabular}
\end{table}

Calibrating on C4 instead of WikiText-2 produces $>$96\% identical head assignments across all three architectures, indicating that per-head sensitivity ranking is an architectural property rather than a dataset-dependent artifact.
The cross-domain PPL gap is minimal for Llama (0.20) and Mistral (0.03), though larger for Qwen (9.27), reflecting Qwen's higher sensitivity to calibration domain.

\subsection{Multi-Seed Stability}

\begin{table}[h]
\centering
\caption{Multi-seed PPL stability (Qwen2.5-7B, $k{=}8$, 3 seeds). Bit distribution is consistent across seeds.}
\label{tab:multi_seed}
\begin{tabular}{@{}lrl@{}}
\toprule
Seed & PPL & Bit Distribution \\
\midrule
0 & 11.02 & $\{2{:}8, 4{:}100, 8{:}4\}$ \\
1 & 10.90 & $\{2{:}8, 4{:}100, 8{:}4\}$ \\
2 & 11.05 & $\{2{:}8, 4{:}100, 8{:}4\}$ \\
\midrule
Mean $\pm$ std & $10.99 \pm 0.065$ & consistent \\
\bottomrule
\end{tabular}
\end{table}

The 95\% confidence interval $[10.86, 11.12]$ is well below uniform 4-bit (15.45), confirming that the improvement is statistically robust.
Bit distribution is identical across all seeds, and 90.2\% of head allocations agree across all three seeds.


%% ============================================================
%% 6  WHEN DOES NON-UNIFORM ALLOCATION HELP?
%% ============================================================
\section{When Does Non-Uniform Allocation Help?}
\label{sec:when_helps}

\subsection{Cross-Architecture Downstream Evaluation}

Table~\ref{tab:downstream} presents downstream task results across 5 architectures under $4\times$ compression.
Only Qwen2.5-7B (spread sensitivity pattern) shows consistent downstream gains (+1.25 to +4.00pp).
The remaining 4 architectures show minor degradation: Yi-1.5 $-0.77$pp, Llama $-0.99$pp, Mistral $-0.48$pp, Gemma-2 $+0.14$pp (noise).

\begin{table}[h]
\centering
\caption{Downstream evaluation ($4\times$ compression, 0-shot). MCKP benefits only Qwen2.5-7B (spread pattern). ARC-Challenge shows the largest per-task sensitivity to allocation strategy across all non-Spread models.}
\label{tab:downstream}
\begin{tabular}{@{}llrrrrl@{}}
\toprule
Model & Pattern & ARC-C $\Delta$ & HellaSwag $\Delta$ & PIQA $\Delta$ & WinoGrande $\Delta$ & Avg $\Delta$ \\
\midrule
Qwen2.5-7B & Spread & +3.16pp & +4.00pp & +1.25pp & +1.97pp & +2.60pp \\
Yi-1.5-9B & Spike & $-$3.24pp & $-$0.60pp & $-$0.33pp & +1.27pp & $-$0.77pp \\
Llama-3.1-8B & Spike & $-$2.65pp & $-$0.14pp & $-$0.65pp & $-$0.15pp & $-$0.99pp \\
Mistral-7B-v0.3 & Flat & $-$1.19pp & --- & --- & --- & $-$0.48pp \\
Gemma-2-9B & Flat & $-$0.17pp & --- & --- & --- & +0.14pp \\
\bottomrule
\end{tabular}
\end{table}

ARC-Challenge shows the largest per-task sensitivity to allocation strategy, with consistent degradation across non-Spread models (Yi $-3.24$pp, Llama $-2.65$pp, Mistral $-1.19$pp).
This is consistent with ARC-Challenge's reliance on multi-step reasoning, where subtle KV cache quantization errors accumulate.

\subsection{Sensitivity Distribution Taxonomy}

We identify three sensitivity distribution shapes that predict allocation benefit:

\begin{itemize}[leftmargin=*,itemsep=2pt]
\item \textbf{Spread} (Qwen2.5-7B): Multiple heads with cross-bitwidth sensitivity gradients. L0 cluster of 4 heads with $26\times$ dominance ratio. Concentration$=0.61$, $t_1/t_5=0.57$. MCKP benefits: $+27.3\%$ PPL, $+2.6$pp downstream.

\item \textbf{Spike} (Yi-1.5-9B, Llama-3.1-8B): Single extreme outlier dominates. Yi L4H3 accounts for 64\% of total delta. Solver assigns 8-bit to the outlier but cannot differentiate remaining heads. Concentration high (Yi $0.71$) but $t_1/t_5$ also high ($0.90$). No MCKP benefit.

\item \textbf{Flat} (Mistral-7B-v0.3, Gemma-2-9B): Near-homogeneous sensitivity. Profiler reports near-zero inter-head differences (Mistral concentration$=0.22$, Gemma-2 $0.23$). Every 2-bit head hurts equally; non-uniform allocation is futile.
\end{itemize}

\paragraph{The spike--spread distinction.}
Yi-1.5 (GQA 8:1, concentration$=0.71$) and Qwen (GQA 7:1, concentration$=0.61$) illustrate why GQA ratio and concentration alone are insufficient predictors.
Yi has higher GQA and higher concentration than Qwen, yet shows no benefit.
The critical difference is distribution \emph{shape}: Qwen's L0 \emph{cluster} of sensitive heads creates optimization opportunities, while Yi's single extreme outlier cannot be exploited.

\subsection{Diagnostic Validation}

The 2-metric diagnostic (concentration $\gtrsim 0.50$ AND $t_1/t_5 \lesssim 0.70$) correctly classifies all 9 tested models:
\begin{itemize}[leftmargin=*,itemsep=1pt]
\item Qwen2.5-7B: $c{=}0.61$, $t_1/t_5{=}0.57$ $\to$ spread $\to$ benefit predicted $\to$ \textbf{correct} (+27.3\%)
\item Yi-1.5-9B: $c{=}0.71$, $t_1/t_5{=}0.90$ $\to$ spike $\to$ no benefit predicted $\to$ \textbf{correct} ($\approx 0\%$)
\item Llama-3.1-8B: $c{=}0.57$, $t_1/t_5{=}0.85$ $\to$ spike $\to$ no benefit predicted $\to$ \textbf{correct} (+0.2\%)
\item Mistral-7B-v0.3: $c{=}0.22$, $t_1/t_5{=}$--- $\to$ flat $\to$ no benefit predicted $\to$ \textbf{correct} ($-0.5\%$)
\item Gemma-2-9B: $c{=}0.23$, $t_1/t_5{=}0.38$ $\to$ flat $\to$ no benefit predicted $\to$ \textbf{correct} ($-0.3\%$)
\end{itemize}
We emphasize that $n{=}9$ provides a consistent pattern, not a statistical claim.

%% ============================================================
%% 7  RUNTIME ANALYSIS
%% ============================================================
\section{Runtime Analysis}
\label{sec:runtime}

\begin{table}[h]
\centering
\caption{Runtime benchmark (Qwen2.5-7B, seq\_len$=$2048). Mixed-precision incurs 8.5\% throughput overhead vs uniform 4-bit due to per-head dequantization.}
\label{tab:runtime}
\begin{tabular}{@{}lrrrr@{}}
\toprule
Method & Prefill (ms) & Decode (ms) & Throughput (tok/s) & Memory (MB) \\
\midrule
BF16 (no quant.) & 45.86 & 33.78 & 29.6 & 21 \\
Uniform 4-bit & 54.25 & 42.46 & 23.6 & 5.25 \\
Constrained MCKP & 57.03 & 46.09 & 21.8 & 5.25 \\
\bottomrule
\end{tabular}
\end{table}

Mixed-precision MCKP incurs an 8.5\% throughput reduction versus uniform 4-bit (21.8 vs 23.6 tok/s), stemming from per-head dequantization kernel overhead.
The MCKP solver itself is negligible (0.005s).
This overhead is an engineering limitation of the current implementation: a fused mixed-precision dequantization kernel could eliminate it.
The memory footprint is identical ($4\times$ compression) since the average bit-width remains the same.

For memory-bound scenarios (long context, large batch sizes), the $4\times$ memory reduction enables deployments that BF16 cannot support.
On 7/8 architectures where non-uniform allocation provides no PPL benefit, the 8.5\% throughput cost is pure overhead; the diagnostic (Section~\ref{sec:when_helps}) helps practitioners avoid this cost by recommending uniform 4-bit when appropriate.

%% ============================================================
%% 8  CONCLUSION
%% ============================================================
\section{Conclusion}
\label{sec:conclusion}

We present an empirical analysis of per-head KV cache quantization on hardware-aligned bit-widths $\{2,4,8,16\}$, identifying three findings:

\textbf{(1)} 2-bit quantization distortions compound non-additively, causing all tested unconstrained solvers to underperform uniform 4-bit across 8 architectures.
This compounding is superlinear after $k{\approx}16$ simultaneous 2-bit heads and exponential after $k{\approx}48$.

\textbf{(2)} A cardinality constraint ($\le k$ heads at 2-bit, $k{\le}8$) prevents compounding across all tested architectures and makes solver choice irrelevant (greedy$=$MCKP, $\Delta{=}0.0000$).
Practitioners need only add one integer parameter to existing pipelines.

\textbf{(3)} Under moderate compression, only 1/8 architectures benefits from non-uniform allocation.
A 2-metric diagnostic (9/9 correct) tells practitioners when to use profiling-based allocation versus uniform 4-bit.

\paragraph{Limitations.}
Our experiments are limited to 7B--14B scale models; scaling behavior to 70B+ models remains an open question (our 72B profiling attempt was terminated due to resource constraints).
Multi-seed confidence intervals are reported for PPL (std$=$0.065) but not for all downstream tasks.
The 8.5\% throughput overhead of the current mixed-precision implementation could be eliminated with optimized kernels but represents a practical concern.
The diagnostic threshold values ($c{\gtrsim}0.50$, $t_1/t_5{\lesssim}0.70$) are derived from $n{=}9$ models and should be treated as an empirical observation, not a statistically validated rule.

\bibliography{references}
\bibliographystyle{plain}

\begin{thebibliography}{10}
\bibitem{hooper2024kvquant}
Coleman Hooper, Sehoon Kim, Haonan Mohber, et al.
\newblock KVQuant: Towards 10 Million Context Length LLM Inference with KV Cache Quantization.
\newblock \emph{arXiv preprint arXiv:2401.18079}, 2024.

\bibitem{liu2024kivi}
Zirui Liu, Jiayi Yuan, Hongye Jin, et al.
\newblock KIVI: A Tuning-Free Asymmetric 2bit Quantization for KV Cache.
\newblock \emph{arXiv preprint arXiv:2402.02750}, 2024.

\bibitem{dong2024qaq}
Haowen Dong, Jiaqi Zheng, et al.
\newblock QAQ: Quality Adaptive Quantization for LLM KV Cache.
\newblock \emph{arXiv preprint arXiv:2403.04643}, 2024.

\bibitem{ratequant2024}
RateQuant Authors.
\newblock RateQuant: Rate-Distortion Optimal KV Cache Quantization.
\newblock \emph{arXiv preprint}, 2024.

\bibitem{nagel2020up}
Markus Nagel, Rana Ali Amjad, Mart van Baalen, et al.
\newblock Up or Down? Adaptive Rounding for Post-Training Quantization.
\newblock \emph{ICML}, 2020.
\end{thebibliography}

\end{document}
